% GENERATED FILE. Edit canonical lessons, then run npm run generate:books.
% canonical-source-hash: fnv1a64:497b6e6dbea607d1
% canonical-lessons: KA-C455-pravasodyama, KA-C455-gamyasthana, KA-C455-buk-madu, KA-C455-margadarsi, KA-C455-preksaniya

\chapter{Planning a Trip}
\label{ch:ka-planning-a-trip}
\hlchaptermodality{455}

\begin{chapteropening}
\textbf{By the end of this chapter:} \emph{I can say tourism, a destination, to make a booking, a guide and worth seeing.}

Complete the last lesson of chapter 455: I can say tourism, a destination, to make a booking, a guide and worth seeing.
\end{chapteropening}

\section[\texorpdfstring{\kn{ಪ್ರವಾಸೋದ್ಯಮ} (pravāsōdyama)}{ಪ್ರವಾಸೋದ್ಯಮ (pravāsōdyama)}]{\kn{ಪ್ರವಾಸೋದ್ಯಮ} (pravāsōdyama) — tourism; the tourist trade}

\label{lesson:KA-C455-pravasodyama}



\begin{quote}
Before the new one: say the Kannada for high tide, then the Kannada for a whirlpool.
\end{quote}

\subsection*{You'll want to know: \kn{ಪ್ರವಾಸೋದ್ಯಮ}}
\textbf{\kn{ಪ್ರವಾಸೋದ್ಯಮ}} — \emph{pravāsōdyama} — \textquotedblleft{}tourism; the tourist trade\textquotedblright{}.

\begin{quote}
\kn{ಈ} \kn{ಜಿಲ್ಲೆಯಲ್ಲಿ} \kn{ಪ್ರವಾಸೋದ್ಯಮ} \kn{ತುಂಬಾ} \kn{ಬೆಳೆದಿದೆ}.
\end{quote}

\begin{quote}
\emph{ī jilleyalli pravāsōdyama tumbā beḷedide.}
\end{quote}

\begin{quote}
Tourism has grown a lot in this district.
\end{quote}

\kn{ಪ್ರವಾಸ}, a trip, plus \kn{ಉದ್ಯಮ}, an industry; \kn{ಪ್ರವಾಸಿ} is a tourist.

\begin{cousinweb}
\textbf{\kn{ಪ್ರವಾಸೋದ್ಯಮ}} (\emph{pravāsōdyama}) is a modern compound built from Sanskrit words: \emph{pravāsa}, \textquotedblleft{}dwelling abroad, a journey\textquotedblright{} (CDIAL 8802), and \emph{udyama}, \textquotedblleft{}an undertaking, effort\textquotedblright{} (CDIAL 2048).
\end{cousinweb}

\begin{grammarlens}[title={What you've built}]
One more word for getting around.
\end{grammarlens}

\subsection*{Your turn}
\begin{itemize}
\raggedright
  \item \emph{Say it:} \emph{pravāsōdyama}
  \item \emph{Say it:} \emph{pravāsōdyama}, once more
  \item \emph{Say it:} \emph{ī jilleyalli pravāsōdyama tumbā beḷedide.} — Tourism has grown a lot in this district.
  \item \emph{Recall:} say the Kannada for high tide, then the Kannada for a whirlpool, then say \emph{pravāsōdyama} again
\end{itemize}

\begin{tcolorbox}[breakable,colback=teal!4,colframe=teal!35!black,title={Before you move on}]
What does \kn{ಪ್ರವಾಸೋದ್ಯಮ} mean? (\textquotedblleft{}Tourism; the tourist trade\textquotedblright{}.) Say it once more.
\end{tcolorbox}
\section[\texorpdfstring{\kn{ಗಮ್ಯಸ್ಥಾನ} (gamyasthāna)}{ಗಮ್ಯಸ್ಥಾನ (gamyasthāna)}]{\kn{ಗಮ್ಯಸ್ಥಾನ} (gamyasthāna) — a destination}

\label{lesson:KA-C455-gamyasthana}



\begin{quote}
Before the new one: say the Kannada for a whirlpool, then the Kannada for tourism.
\end{quote}

\subsection*{You'll want to know: \kn{ಗಮ್ಯಸ್ಥಾನ}}
\textbf{\kn{ಗಮ್ಯಸ್ಥಾನ}} — \emph{gamyasthāna} — \textquotedblleft{}a destination\textquotedblright{}.

\begin{quote}
\kn{ನಮ್ಮ} \kn{ಬಸ್} \kn{ಗಮ್ಯಸ್ಥಾನವನ್ನು} \kn{ಮಧ್ಯರಾತ್ರಿ} \kn{ತಲುಪಿತು}.
\end{quote}

\begin{quote}
\emph{namma bas gamyasthānavannu madhyarātri talupitu.}
\end{quote}

\begin{quote}
Our bus reached its destination at midnight.
\end{quote}

\kn{ಗಮ್ಯ}, to be gone to, plus \kn{ಸ್ಥಾನ}, a place.

\begin{cousinweb}
\textbf{\kn{ಗಮ್ಯಸ್ಥಾನ}} (\emph{gamyasthāna}) is built from two Sanskrit words: \emph{gamya}, \textquotedblleft{}to be gone to, accessible\textquotedblright{} (CDIAL 4032), and \emph{sthāna}, \textquotedblleft{}a place\textquotedblright{} (CDIAL 13753).
\end{cousinweb}

\begin{grammarlens}[title={What you've built}]
One more word for getting around.
\end{grammarlens}

\subsection*{Your turn}
\begin{itemize}
\raggedright
  \item \emph{Say it:} \emph{gamyasthāna}
  \item \emph{Say it:} \emph{gamyasthāna}, once more
  \item \emph{Say it:} \emph{namma bas gamyasthānavannu madhyarātri talupitu.} — Our bus reached its destination at midnight.
  \item \emph{Recall:} say the Kannada for a whirlpool, then the Kannada for tourism, then say \emph{gamyasthāna} again
\end{itemize}

\begin{tcolorbox}[breakable,colback=teal!4,colframe=teal!35!black,title={Before you move on}]
What does \kn{ಗಮ್ಯಸ್ಥಾನ} mean? (\textquotedblleft{}A destination\textquotedblright{}.) Say it once more.
\end{tcolorbox}
\section[\texorpdfstring{\kn{ಬುಕ್} \kn{ಮಾಡು} (buk māḍu)}{ಬುಕ್ ಮಾಡು (buk māḍu)}]{\kn{ಬುಕ್} \kn{ಮಾಡು} (buk māḍu) — to make a booking; to book}

\label{lesson:KA-C455-buk-madu}



\begin{quote}
Before the new one: say the Kannada for tourism, then the Kannada for a destination.
\end{quote}

\subsection*{You'll want to know: \kn{ಬುಕ್} \kn{ಮಾಡು}}
\textbf{\kn{ಬುಕ್} \kn{ಮಾಡು}} — \emph{buk māḍu} — \textquotedblleft{}to make a booking; to book\textquotedblright{}.

\begin{quote}
\kn{ನಾನು} \kn{ಆಪ್ನಲ್ಲಿ} \kn{ಎರಡು} \kn{ಟಿಕೆಟ್} \kn{ಬುಕ್} \kn{ಮಾಡಿದೆ}.
\end{quote}

\begin{quote}
\emph{nānu āpnalli eraḍu ṭikeṭ buk māḍide.}
\end{quote}

\begin{quote}
I booked two tickets on an app.
\end{quote}

Past (I …): \kn{ಬುಕ್} \kn{ಮಾಡಿದೆ}, \emph{buk māḍide}.

Everyday speech for \kn{ಕಾಯ್ದಿರಿಸು}, to reserve.

\begin{cousinweb}
\textbf{\kn{ಬುಕ್} \kn{ಮಾಡು}} (\emph{buk māḍu}) joins a loan from English \emph{book}, \textquotedblleft{}to reserve\textquotedblright{}, and the native verb \kn{ಮಾಡು}, to do.
\end{cousinweb}

\begin{grammarlens}[title={What you've built}]
One more word for getting around.
\end{grammarlens}

\subsection*{Your turn}
\begin{itemize}
\raggedright
  \item \emph{Say it:} \emph{buk māḍu}
  \item \emph{Say it:} \emph{buk māḍu}, once more
  \item \emph{Say it:} \emph{nānu āpnalli eraḍu ṭikeṭ buk māḍide.} — I booked two tickets on an app.
  \item \emph{Recall:} say the Kannada for tourism, then the Kannada for a destination, then say \emph{buk māḍu} again
\end{itemize}

\begin{tcolorbox}[breakable,colback=teal!4,colframe=teal!35!black,title={Before you move on}]
What does \kn{ಬುಕ್} \kn{ಮಾಡು} mean? (\textquotedblleft{}To make a booking; to book\textquotedblright{}.) Say it once more.
\end{tcolorbox}
\section[\texorpdfstring{\kn{ಮಾರ್ಗದರ್ಶಿ} (mārgadarśi)}{ಮಾರ್ಗದರ್ಶಿ (mārgadarśi)}]{\kn{ಮಾರ್ಗದರ್ಶಿ} (mārgadarśi) — a guide}

\label{lesson:KA-C455-margadarsi}



\begin{quote}
Before the new one: say the Kannada for a destination, then the Kannada for to make a booking.
\end{quote}

\subsection*{You'll want to know: \kn{ಮಾರ್ಗದರ್ಶಿ}}
\textbf{\kn{ಮಾರ್ಗದರ್ಶಿ}} — \emph{mārgadarśi} — \textquotedblleft{}a guide\textquotedblright{}.

\begin{quote}
\kn{ಮಾರ್ಗದರ್ಶಿ} \kn{ನಮಗೆ} \kn{ಕೋಟೆಯ} \kn{ಇತಿಹಾಸ} \kn{ಹೇಳಿದನು}.
\end{quote}

\begin{quote}
\emph{mārgadarśi namage kōṭeya itihāsa hēḷidanu.}
\end{quote}

\begin{quote}
The guide told us the history of the fort.
\end{quote}

Plural: \kn{ಮಾರ್ಗದರ್ಶಿಗಳು}, \emph{mārgadarśigaḷu}.

\kn{ಮಾರ್ಗ}, a way, plus -\kn{ದರ್ಶಿ}, one who shows; a guidebook can be one too.

\begin{cousinweb}
\textbf{\kn{ಮಾರ್ಗದರ್ಶಿ}} (\emph{mārgadarśi}) is built from Sanskrit \emph{mārga}, \textquotedblleft{}a way\textquotedblright{} (CDIAL 10071), and \emph{-darśin}, \textquotedblleft{}one who shows\textquotedblright{}, from \emph{dṛś}, \textquotedblleft{}to see\textquotedblright{}; Sanskrit also has \emph{mārgadarśaka}, \textquotedblleft{}a guide\textquotedblright{} (Monier-Williams).
\end{cousinweb}

\begin{grammarlens}[title={What you've built}]
One more word for getting around.
\end{grammarlens}

\subsection*{Your turn}
\begin{itemize}
\raggedright
  \item \emph{Say it:} \emph{mārgadarśi}
  \item \emph{Say it:} \emph{mārgadarśi}, once more
  \item \emph{Say it:} \emph{mārgadarśi namage kōṭeya itihāsa hēḷidanu.} — The guide told us the history of the fort.
  \item \emph{Recall:} say the Kannada for a destination, then the Kannada for to make a booking, then say \emph{mārgadarśi} again
\end{itemize}

\begin{tcolorbox}[breakable,colback=teal!4,colframe=teal!35!black,title={Before you move on}]
What does \kn{ಮಾರ್ಗದರ್ಶಿ} mean? (\textquotedblleft{}A guide\textquotedblright{}.) Say it once more.
\end{tcolorbox}
\section[\texorpdfstring{\kn{ಪ್ರೇಕ್ಷಣೀಯ} (prēkṣaṇīya)}{ಪ್ರೇಕ್ಷಣೀಯ (prēkṣaṇīya)}]{\kn{ಪ್ರೇಕ್ಷಣೀಯ} (prēkṣaṇīya) — worth seeing}

\label{lesson:KA-C455-preksaniya}



\begin{quote}
Before the new one: say the Kannada for to make a booking, then the Kannada for a guide.
\end{quote}

\subsection*{You'll want to know: \kn{ಪ್ರೇಕ್ಷಣೀಯ}}
\textbf{\kn{ಪ್ರೇಕ್ಷಣೀಯ}} — \emph{prēkṣaṇīya} — \textquotedblleft{}worth seeing\textquotedblright{}.

\begin{quote}
\kn{ಈ} \kn{ಊರಿನಲ್ಲಿ} \kn{ತುಂಬಾ} \kn{ಪ್ರೇಕ್ಷಣೀಯ} \kn{ದೇವಸ್ಥಾನಗಳಿವೆ}.
\end{quote}

\begin{quote}
\emph{ī ūrinalli tumbā prēkṣaṇīya dēvasthānagaḷive.}
\end{quote}

\begin{quote}
There are many temples worth seeing in this town.
\end{quote}

\kn{ಪ್ರೇಕ್ಷಕ}, a viewer, comes from the same root.

\begin{cousinweb}
\textbf{\kn{ಪ್ರೇಕ್ಷಣೀಯ}} (\emph{prēkṣaṇīya}) is borrowed from Sanskrit \emph{prēkṣaṇīya}, \textquotedblleft{}to be seen, visible\textquotedblright{} (Monier-Williams), on \emph{prēkṣaṇa}, \textquotedblleft{}viewing, a show\textquotedblright{} (CDIAL 8991).
\end{cousinweb}

\begin{grammarlens}[title={What you've built}]
One more word for getting around.
\end{grammarlens}

\subsection*{Your turn}
\begin{itemize}
\raggedright
  \item \emph{Say it:} \emph{prēkṣaṇīya}
  \item \emph{Say it:} \emph{prēkṣaṇīya}, once more
  \item \emph{Say it:} \emph{ī ūrinalli tumbā prēkṣaṇīya dēvasthānagaḷive.} — There are many temples worth seeing in this town.
  \item \emph{Recall:} say the Kannada for to make a booking, then the Kannada for a guide, then say \emph{prēkṣaṇīya} again
\end{itemize}

\begin{tcolorbox}[breakable,colback=teal!4,colframe=teal!35!black,title={Before you move on}]
What does \kn{ಪ್ರೇಕ್ಷಣೀಯ} mean? (\textquotedblleft{}Worth seeing\textquotedblright{}.) Say it once more.
\end{tcolorbox}
